\documentclass[12pt, reqno]{amsart}
\usepackage{geometry}
\geometry{margin=1in}
\usepackage{amsmath, amssymb, amsthm}
\usepackage{mathrsfs, mathtools, stmaryrd}
\usepackage[colorlinks=true, citecolor=blue, linkcolor=red]{hyperref}
\usepackage{tikz-cd}
\usepackage{enumitem}
\usepackage{bbm}
\usepackage{bm}

% --- Theorem Environments ---
\theoremstyle{plain}
\newtheorem{theorem}{Theorem}[section]
\newtheorem{lemma}[theorem]{Lemma}
\newtheorem{proposition}[theorem]{Proposition}
\newtheorem{corollary}[theorem]{Corollary}
\newtheorem{axiom}{Axiom}[section]

\theoremstyle{definition}
\newtheorem{definition}[theorem]{Definition}
\newtheorem{example}[theorem]{Example}
\newtheorem{construction}[theorem]{Construction}

\theoremstyle{remark}
\newtheorem{remark}[theorem]{Remark}
\newtheorem{convention}[theorem]{Convention}

% --- Special Operators ---
\DeclareMathOperator{\id}{id}
\DeclareMathOperator{\tr}{tr}
\DeclareMathOperator{\Tr}{Tr}
\DeclareMathOperator{\Spec}{Spec}
\DeclareMathOperator{\Ker}{Ker}
\DeclareMathOperator{\rank}{rank}
\DeclareMathOperator{\Hom}{Hom}
\DeclareMathOperator{\ad}{ad}
\DeclareMathOperator{\Exp}{Exp}
\DeclareMathOperator{\Hol}{Hol}

% --- Shortcuts and Symbols ---
\newcommand{\cA}{\mathcal{A}}
\newcommand{\cO}{\mathcal{O}}
\newcommand{\cH}{\mathcal{H}}
\newcommand{\cK}{\mathcal{K}}
\newcommand{\cL}{\mathcal{L}}
\newcommand{\cN}{\mathcal{N}}
\newcommand{\cP}{\mathcal{P}}
\newcommand{\bC}{\mathbb{C}}
\newcommand{\bR}{\mathbb{R}}
\newcommand{\bZ}{\mathbb{Z}}
\newcommand{\bN}{\mathbb{N}}

% --- Lambda Framework ---
\newcommand{\lstar}{\lambda_{*}}
\newcommand{\Lam}{\boldsymbol{\Lambda}}
\newcommand{\lam}{\lambda}
\newcommand{\val}{v_{\lambda}}
\newcommand{\unit}{U}
\newcommand{\seam}{\Gamma}

% --- Operator Symbols ---
\newcommand{\bindu}{\odot}
\newcommand{\seamop}{\mathcal{S}_{\seam}}
\newcommand{\release}{\mathcal{R}}
\newcommand{\aghora}{\mathcal{A}}
\newcommand{\entropy}{\mathcal{E}}
\newcommand{\bond}{B}
\newcommand{\shift}{L}
\newcommand{\scale}{S}
\newcommand{\fred}{\widehat{K}}
\newcommand{\eco}{C_{v}}
\newcommand{\ecoflow}{e^{tD_{v}}}

% --- Four Ways ---
\newcommand{\wayone}{\mathsf{W}_{1}}
\newcommand{\waytwo}{\mathsf{W}_{2}}
\newcommand{\waythree}{\mathsf{W}_{3}}
\newcommand{\wayfour}{\mathsf{W}_{4}}

% --- Sanskrit Mantra Operators ---
\newcommand{\OM}{\mathsf{A}_{OM}}
\newcommand{\NA}{\mathsf{N}}
\newcommand{\MA}{\mathsf{M}}
\newcommand{\SHI}{\mathsf{Sh}}
\newcommand{\VA}{\mathsf{V}}
\newcommand{\YA}{\mathsf{Y}}

\title{A Unified Framework for $\lambda$-Structures and Operator Algebras}
\author{Four-Way Representation Theory}
\date{\today}

\begin{document}

\maketitle

\begin{abstract}
This work presents a complete mathematical framework organizing structures into four canonical representations: 
Way-1 (intrinsic multidimensional numbers), Way-2 (relational/response numbers), Way-3 (reference-scale 
normalization), and Way-4 (operator evolution). We develop the foundational $\lambda$-valued field, construct 
the operator algebras that emerge from it, and demonstrate how all four representations are translations of a 
single underlying system. Applications to spectral theory, Fredholm operators, and symbolic mantra algebras are 
developed in full detail.
\end{abstract}

\tableofcontents

\newpage

\section{Introduction: The Four-Way Framework}

This work unifies two parallel developments: (1) the $\lambda$-number theory where numbers are constructed 
from a universal reference constant $\lstar$, and (2) the operator-theoretic framework where systems evolve 
through exponential change operators. The synthesis is achieved through a four-fold representation:

\begin{description}
\item[Way-1 (\wayone):] Intrinsic multidimensional structure. Objects are vectors
    $\mathbf{x} = (x_1,\ldots,x_n)$ with components in a $\lambda$-valued field.
    
\item[Way-2 (\waytwo):] Relational/response numbers. Objects defined through structural relations
    $\lambda = dX/dY$ and transformation laws.
    
\item[Way-3 (\waythree):] Reference-scale normalization. All quantities expressed as multiples of a universal
    constant $\lstar$: $x = n\lstar$ or $x = u\lstar^{k}$.
    
\item[Way-4 (\wayfour):] Operator evolution. Systems described by transformation operators
    $\mid \psi_{\text{final}}\rangle = U \mid \psi_0\rangle$ and exponential flows $e^{tD}$.
\end{description}

These four ways are not independent theories but complementary representations of the same underlying 
mathematical reality, each revealing different aspects of structure, relation, scale, and dynamics.

\section{The $\lambda$-Valued Field}

\begin{definition}[Reference Constant]
Let $\lstar \neq 0$ be a distinguished constant called the \textbf{universal reference constant}.
\end{definition}

\begin{definition}[$\lambda$-Valued Field]
Define
\[
K = \{ \lstar^{k}u \mid k \in \bZ,\; u \in \unit \}
\]
where $\unit$ is a compact multiplicative group of \textbf{units} satisfying $\val(u)=0$. The 
\textbf{$\lambda$-valuation} is
\[
\val(\lstar^{k}u) = k.
\]
\end{definition}

\begin{proposition}
$K$ is a field with valuation $\val: K^{\times} \to \bZ$ satisfying
\begin{align*}
\val(xy) &= \val(x) + \val(y), \\
\val(x+y) &\geq \min\{\val(x),\val(y)\}.
\end{align*}
Moreover, $K^{\times} \cong \bZ \times \unit$.
\end{proposition}

\section{Way-1: Intrinsic Multidimensional Numbers}

\begin{definition}[Way-1 Number Space]
For $n \in \bN$, define the \textbf{Way-1 number space} as
\[
\cN_{n} = K^{n} = \{ \mathbf{x} = (x_1,\ldots,x_n) \mid x_i \in K \}.
\]
\end{definition}

\begin{definition}[Coordinate Algebra]
Define coordinatewise operations:
\begin{align*}
\mathbf{x} + \mathbf{y} &= (x_1 + y_1,\ldots,x_n + y_n), \\
\alpha\mathbf{x} &= (\alpha x_1,\ldots,\alpha x_n), \\
\mathbf{x} \odot \mathbf{y} &= (x_1y_1,\ldots,x_ny_n).
\end{align*}
The triple $(\cN_n, +, \odot)$ is a commutative $K$-algebra.
\end{definition}

\begin{theorem}[Basis and Decomposition]
The standard basis vectors $e_i = (0,\ldots,1,\ldots,0)$ satisfy:
\begin{align*}
e_i \odot e_i &= e_i, \\
e_i \odot e_j &= 0 \quad (i \neq j), \\
\mathbf{x} &= \sum_{i=1}^{n} x_i e_i.
\end{align*}
Thus $\cN_n \cong Ke_1 \oplus \cdots \oplus Ke_n$ as algebras.
\end{theorem}

\begin{definition}[Valuation on Way-1]
For $\mathbf{x} = (x_1,\ldots,x_n)$, define
\[
\val(\mathbf{x}) = \min_{1\leq i\leq n} \val(x_i).
\]
Then $\val(\mathbf{x} \odot \mathbf{y}) = \val(\mathbf{x}) + \val(\mathbf{y})$ when the minima align.
\end{definition}

\begin{theorem}[Graded Structure]
\[
\cN_n = \bigcup_{k\in\bZ} \cN_n^{(k)}, \quad \cN_n^{(k)} = \{ \mathbf{x} \mid \val(\mathbf{x}) = k \},
\]
and $\cN_n^{(k)} \odot \cN_n^{(m)} \subseteq \cN_n^{(k+m)}$. Hence $\cN_n$ is a graded algebra.
\end{theorem}

\section{Way-2: Relational/Response Numbers}

\begin{definition}[Structural Response]
A \textbf{response number} is a quantity $\lambda$ defined through a relation
\[
\lambda = \frac{dX}{dY}
\]
or more generally $\lambda_{ij} = \partial x_i / \partial x_j$. These are Jacobian-type numbers
measuring relative change between structures.
\end{definition}

\begin{definition}[Relational Algebra]
A \textbf{relational algebra} is a set equipped with:
\begin{itemize}
    \item A collection of response numbers $\{\lambda_{\alpha}\}$
    \item Transformation laws expressing how structures relate
    \item Invariant relations preserved under transformations
\end{itemize}
\end{definition}

\begin{example}[Differential Relations]
For observables $F,G$, the response relation
\[
\frac{dF}{dt} = \lambda \frac{dG}{dt}
\]
defines $\lambda$ as a structural coupling constant.
\end{example}

\begin{definition}[Seam Invariant]
A \textbf{seam} is an invariant relation
\[
\seam = \frac{\lambda_p}{\lambda_v}
\]
that must be preserved under all admissible transformations.
\end{definition}

\begin{definition}[Seam Operator]
For quantities $A,B$ with a seam $\seam$, define
\[
\seamop(A,B) = \lim_{\tau \to \tau_0 \text{ along seam}} \frac{A(\tau)}{B(\tau)}
\]
where the limit is taken while preserving $A/B = \seam$.
\end{definition}

\section{Way-3: Reference-Scale Normalization}

\begin{definition}[Normalized Quantities]
Every quantity $x \in K$ can be written uniquely as
\[
x = n\lstar \quad \text{or} \quad x = u\lstar^{k}
\]
with $n$ dimensionless and $u \in \unit$. The integer $k = \val(x)$ is the \textbf{scale order}.
\end{definition}

\begin{definition}[Harmonic Representation]
A system in Way-3 is described by dimensionless coefficients $\{n_{\alpha}\}$ such that
\[
X_{\alpha} = n_{\alpha}\lstar^{\,k_{\alpha}}.
\]
Transformations preserve the reference scale $\lstar$ while acting on the coefficients.
\end{definition}

\begin{example}[Fredholm in Way-3]
The Fredholm equation becomes
\[
n(\lambda) = n_0(\lambda) + \int \widetilde{K}(\lambda,\lambda')\, n(\lambda')\, d\lambda'
\]
with $\widetilde{K} = K/\lstar$ dimensionless.
\end{example}

\section{Way-4: Operator Evolution}

\begin{definition}[Operator Algebra]
A \textbf{Way-4 system} consists of:
\begin{itemize}
    \item A Hilbert space $\cH$ of states $\mid \psi\rangle$
    \item A collection of operators $\{O_i\}$ acting on $\cH$
    \item Composition rules $O_i \circ O_j$
    \item Exponential flows $e^{tD}$ generated by derivations $D$
\end{itemize}
\end{definition}

\begin{definition}[Exponential Change Operator]
For a generator $D_v$, define
\[
\eco(t) = \ecoflow = \sum_{k=0}^{\infty} \frac{t^k}{k!} D_v^{\,k}.
\]
This satisfies the semigroup law $\eco(t+s) = \eco(t)\eco(s)$ and the differential equation
\[
\frac{d}{dt}\eco(t) = D_v\eco(t).
\]
\end{definition}

\begin{theorem}[Broken Symmetry Decay]
If $U$ satisfies $D_v U + \lambda U = 0$, then under the flow
\[
U(\phi_t(p)) = U(p)e^{-\lambda t}.
\]
Thus $\lambda$ is the decay rate associated to broken symmetry.
\end{theorem}

\begin{definition}[Fredholm Operator]
The $\lambda$-Fredholm operator is
\[
(I - \fred_{\lambda})\psi = f
\]
with kernel $K(\lambda,\lambda')$ encoding system interactions. Solvability occurs when
\[
\det(I - \fred_{\lambda}) = 0.
\]
\end{definition}

\begin{definition}[Winding-Mellin Operator]
The full scale-evolution operator is
\[
(Wf)(t,\theta) = -\int K(\Delta t;\theta,\theta')\, e^{\sigma\Delta t} f(t-\Delta t,\theta')\, d\theta' d\Delta t
+ \nu\Delta_{\theta}f + \mu\partial_t f + \kappa J_{\theta}f
\]
where:
\begin{itemize}
    \item $t = \log r$ is logarithmic scale
    \item $\theta$ is angular/phase coordinate
    \item $K$ encodes memory across scales
    \item $\nu\Delta_{\theta}$ is phase diffusion
    \item $\mu\partial_t$ is scale drift
    \item $\kappa J_{\theta}$ is winding current
\end{itemize}
\end{definition}

\section{Canonical Operators and Their Relations}

\subsection{Shift and Scale}

\begin{definition}[Hilbert Space Completion]
Let $\cH = \ell^2(K)$ with orthonormal basis $\{e_i\}_{i=1}^{\infty}$.
\end{definition}

\begin{definition}[Scale Operator]
\[
\scale e_i = i e_i, \quad \scale(x_1,x_2,\ldots) = (x_1,2x_2,3x_3,\ldots)
\]
with domain $\{ x \mid \sum i^2 |x_i|^2 < \infty \}$.
\end{definition}

\begin{definition}[Shift Operator]
\[
\shift e_i = e_{i+1}, \quad \shift(x_1,x_2,\ldots) = (0,x_1,x_2,\ldots)
\]
$\shift$ is an isometry: $\|\shift x\| = \|x\|$.
\end{definition}

\begin{theorem}[Fundamental Commutation]
On common domain,
\[
[\scale,\shift] = \shift, \quad [\scale,\shift^*] = -\shift^*.
\]
Consequences:
\begin{align*}
\shift^*\shift &= I, \quad \shift\shift^* = I - P_1, \\
\shift^{*t}\scale\shift^t &= \scale + tI, \\
\shift\scale &= (\scale + I)\shift.
\end{align*}
\end{theorem}

\begin{theorem}[Fluctuation Growth]
For $x = e_1$, the orbit sum satisfies
\[
\Big\| \sum_{t=0}^{N-1} \shift^t e_1 \Big\|^2 = N,
\]
so $\| \sum_{t=0}^{N-1} \shift^t x \| = O(\sqrt{N})$ for finitely supported $x$.
\end{theorem}

\subsection{Bindu Operator}

\begin{definition}[Bindu Operator]
The \textbf{Bindu operator} $\bindu$ maps magnitudes to phase:
\[
\bindu(x) = e^{i\Phi(x)}
\]
where $\Phi(x)$ is a phase functional. In operator form,
\[
\bindu: \cH \to S^1 \subset \bC.
\]
\end{definition}

\begin{theorem}[Regularization Property]
For indeterminate forms such as $0/0$, $0\cdot\infty$, or $\infty/\infty$, $\bindu$ produces
a well-defined phase:
\[
\bindu(0/0) = e^{i(\Phi(0)-\Phi(0))} = 1.
\]
Thus $\bindu$ regularizes scale singularities.
\end{theorem}

\subsection{Entropic Operator}

\begin{definition}[Entropic Operator]
For a state space $S$ and resolution scale $\lambda$, define
\[
\entropy_{\lambda}: S \to \{ O_i \}_{i=1}^{N(\lambda)}
\]
where $N(\lambda)$ is the number of observable objects at scale $\lambda$, satisfying
$N(\lambda_1) > N(\lambda_2)$ whenever $\lambda_1 < \lambda_2$.
\end{definition}

\begin{proposition}
The object count relates to entropy via
\[
N(\lambda) \sim e^{H(\lambda)}
\]
where $H(\lambda)$ is the Shannon entropy at scale $\lambda$.
\end{proposition}

\subsection{Aghora Operator}

\begin{definition}[Aghora Operator]
Let $\aghora$ be an involution satisfying
\[
\aghora^2 = I, \quad \aghora G \aghora = -G
\]
for a distortion operator $G$. Thus $\aghora$ inverts distortion while preserving the underlying structure.
\end{definition}

\subsection{Bond-Release Operator}

\begin{definition}[Bond Operator]
Let $\bond$ be a projection onto a distinguished subspace $\cH_{\text{bond}} \subset \cH$, representing
attached or constrained states.
\end{definition}

\begin{definition}[Release Semigroup]
Define the \textbf{release operator}
\[
\release(\tau) = e^{-\tau K}, \quad \tau \geq 0
\]
where $K$ is a positive generator. Then
\[
\|\release(\tau)\psi\| \leq \|\psi\|, \quad \release(\tau)\cH_{\text{bond}} \to \cH_{\text{free}} \text{ as } \tau\to\infty.
\end{definition}

\section{The Four-Way Translation Dictionary}

\begin{center}
\begin{tabular}{|c|c|c|c|}
\hline
\textbf{Concept} & \textbf{Way-1} & \textbf{Way-2} & \textbf{Way-3} & \textbf{Way-4} \\
\hline
Number & $(x_1,\ldots,x_n)$ & $\lambda = dX/dY$ & $n\lstar$ & $\mid \psi\rangle$ \\
\hline
Addition & Componentwise & Relational composition & Scale addition & Operator sum \\
\hline
Multiplication & $\odot$ coordinatewise & Response composition & Scale multiplication & Operator composition \\
\hline
Evolution & $X \mapsto MX$ & $dX/dt = \lambda Y$ & $n(t) = n_0 e^{-\lambda t}$ & $e^{tD}\mid\psi_0\rangle$ \\
\hline
Invariant & Basis idempotents & Seam $\seam = \lambda_p/\lambda_v$ & Reference $\lstar$ & Spectrum $\Spec(D)$ \\
\hline
\end{tabular}
\end{center}

\section{Sanskrit Mantras in Operator Form}

\subsection{OM as First Eigenmode}

\begin{definition}[OM State]
Let $\mid \Omega\rangle$ be the primordial state satisfying the eigenvalue equation
\[
\Delta \mid \Omega\rangle = -\omega_1 \mid \Omega\rangle
\]
with minimal nonzero frequency $\omega_1 > 0$.
\end{definition}

\begin{definition}[OM Construction]
\[
\mid \Omega\rangle = N\, S_{a,u}(C_a\mid 0\rangle \otimes C_u\mid 0\rangle)
\]
where $C_a$ is creation, $C_u$ is preservation, and $S_{a,u}$ is fusion.
\end{definition}

\subsection{Om Namah Shivaya: Four-Way Representation}

Let the mantra be $\OM \to \NA \to \MA \to \SHI \to \VA \to \YA$, representing:
\begin{itemize}
    \item $\OM$: primordial excitation
    \item $\NA$: negation/grounding
    \item $\MA$: containment/stabilization
    \item $\SHI$: purification/illumination
    \item $\VA$: expansion/flow
    \item $\YA$: union/integration
\end{itemize}

\subsubsection{Way-1: Structural Vector}

\[
\Psi = (s,g,c,p,e,u) \in K^6
\]
with transformation $\Psi' = M_{NS}\Psi$ where
\[
M_{NS} = \begin{pmatrix}
1 & 0 & 0 & 0 & 0 & 0 \\
\alpha & 1 & 0 & 0 & 0 & 0 \\
0 & \beta & 1 & 0 & 0 & 0 \\
0 & 0 & \gamma & 1 & 0 & 0 \\
0 & 0 & 0 & \delta & 1 & 0 \\
0 & 0 & 0 & 0 & \eta & 1
\end{pmatrix}.
\]

\subsubsection{Way-2: Response Chain}

\[
\OM \xrightarrow{\lambda_{GS}} \NA \xrightarrow{\lambda_{CG}} \MA \xrightarrow{\lambda_{PC}} \SHI
\xrightarrow{\lambda_{EP}} \VA \xrightarrow{\lambda_{UE}} \YA
\]
with differential relations:
\[
\frac{dG}{d\tau} = \lambda_{GS}S,\; \frac{dC}{d\tau} = \lambda_{CG}G,\; \frac{dP}{d\tau} = \lambda_{PC}C,\;
\frac{dE}{d\tau} = \lambda_{EP}P,\; \frac{dU}{d\tau} = \lambda_{UE}E.
\]

\subsubsection{Way-3: Harmonic Alignment}

\[
\OM = n_S\lstar,\; \NA = n_G\lstar,\; \MA = n_C\lstar,\; \SHI = n_P\lstar,\; \VA = n_E\lstar,\; \YA = n_U\lstar
\]
with $n_U = (a_5a_4a_3a_2a_1)n_S$.

\subsubsection{Way-4: Operator Evolution}

\[
\mid \Psi_{\text{final}}\rangle = \YA \VA \SHI \MA \NA \OM \mid \Psi_0\rangle.
\]

Fredholm form:
\[
(I - \widehat{K}_{NS})\Psi = \Phi_{\OM}, \quad \widehat{K}_{NS} = \widehat{K}_{\YA} \circ \widehat{K}_{\VA} \circ \widehat{K}_{\SHI} \circ \widehat{K}_{\MA} \circ \widehat{K}_{\NA}.
\]

Resonance condition:
\[
\det(I - \widehat{K}_{NS}) = 0.
\]

\section{Unified Operator Algebra}

\begin{definition}[Five-Operator System]
The complete framework is generated by five fundamental operators:
\begin{align*}
\entropy &: \text{entropy/object generation} \\
\bindu &: \text{phase regularization} \\
\seamop &: \text{invariant relation selection} \\
\aghora &: \text{distortion inversion} \\
\release &: \text{bound relaxation}
\end{align*}
\end{definition}

\begin{theorem}[Unified Evolution]
For a state $\mid \psi\rangle$, the full transformation is
\[
\mid \psi_{\text{final}}\rangle = \release(\tau)\,\aghora\,\bindu\,\seamop\,\bond \mid \psi_0\rangle.
\]
\end{theorem}

\begin{theorem}[Spectral Fluctuation Law]
For the shift-scale pair $(\scale,\shift)$ on $\ell^2(K)$:
\begin{enumerate}
    \item $\scale$ has discrete spectrum $\sigma(\scale) = \{1,2,3,\ldots\}$
    \item $\shift$ is an isometry with $\shift^*\shift = I$, $\shift\shift^* = I - P_1$
    \item $\|\sum_{t=0}^{N-1}\shift^t x\| = O(\sqrt{N})$ for finitely supported $x$
    \item For normalized orbit sums $\Phi_N(x) = \frac{1}{\sqrt{N}}\sum_{t=0}^{N-1}\shift^t x$,
          entropy grows as $H(\Phi_N(e_1)) = \log N$
\end{enumerate}
\end{theorem}

\begin{theorem}[$\lambda$-Fredholm Completeness]
Let $(I - \fred_{\lambda})\psi = f$ be a $\lambda$-Fredholm system with kernel $K_{\lambda}$.
Then:
\begin{enumerate}
    \item Solvability occurs at $\det(I - \fred_{\lambda}) = 0$
    \item Under ECO evolution, $e^{-tD_v}(I - \fred_{\lambda})\psi = f_t$
    \item The winding-Mellin extension incorporates scale memory and phase dynamics
\end{enumerate}
\end{theorem}

\section{Conclusion: The Four-Way Unity}

\begin{theorem}[Representation Completeness]
Every mathematical structure admitting:
\begin{itemize}
    \item intrinsic coordinates (Way-1),
    \item relational laws (Way-2),
    \item scale normalization (Way-3),
    \item and transformational dynamics (Way-4)
\end{itemize}
can be represented in any of the four ways, with translation functors relating them. Conversely,
any system expressed in one way admits equivalent formulations in the other three, revealing the
underlying unity of structural, relational, scalar, and dynamical descriptions.
\end{theorem}

\begin{center}
\fbox{
\begin{minipage}{0.8\textwidth}
\centering
\textbf{The Four-Way Principle}\\
Structure $\leftrightarrow$ Relation $\leftrightarrow$ Scale $\leftrightarrow$ Transformation
\end{minipage}
}
\end{center}

\begin{example}[Mantra Synthesis]
For $\OM \NA \MA \SHI \VA \YA$ (Om Namah Shivaya):
\begin{align*}
\text{Way-1} &: (s,g,c,p,e,u) \mapsto (s',g',c',p',e',u') \\
\text{Way-2} &: \OM \xrightarrow{\lambda} \NA \xrightarrow{\lambda} \MA \xrightarrow{\lambda} \SHI \xrightarrow{\lambda} \VA \xrightarrow{\lambda} \YA \\
\text{Way-3} &: n_S\lstar \to n_G\lstar \to n_C\lstar \to n_P\lstar \to n_E\lstar \to n_U\lstar \\
\text{Way-4} &: \mid \Psi_{\text{final}}\rangle = \YA \VA \SHI \MA \NA \OM \mid \Psi_0\rangle
\end{align*}
All four describe the same transformation: primordial excitation, grounding, stabilization,
purification, expansion, and union.
\end{example}

\begin{thebibliography}{99}

\bibitem{lambda-field} Author. ``The $\lambda$-Valued Field and Its Graded Structure.'' (2026).

\bibitem{four-ways} Author. ``Four Canonical Representations of Mathematical Structure.'' (2026).

\bibitem{fredholm-eco} Author. ``$\lambda$-Fredholm Operators and Exponential Change.'' (2026).

\bibitem{winding-mellin} Author. ``Winding-Mellin Operators and Scale Memory.'' (2026).

\bibitem{mantra-operators} Author. ``Sanskrit Mantras as Operator Algebras.'' (2026).

\end{thebibliography}

\end{document}